% !PRINTED 406
\weberpage{406}
If one sets corresponding coefficients equal to one another,
it follows that
\[
\begin{aligned}
1.&\quad \alpha^2 - 2\beta = 36,\\
2.&\quad \beta^2 + 2\delta - 2\alpha\gamma = 528,\\
3.&\quad \gamma^2 + 2\alpha\varepsilon - 2\beta\delta = 3897,\\
4.&\quad \delta^2 - 2\gamma\varepsilon = 14354,\\
5.&\quad \varepsilon^2 = 21025,
\end{aligned}
\]
and it is now a matter of seeking integral values of $\alpha, \beta, \gamma, \delta, \varepsilon$
that satisfy these five equations. First of all, from 5. one obtains
\[
\varepsilon = \sqrt{21025} = 145 = 5\,.\,29;
\]
and one may take $\varepsilon$ positive, because the term independent of $x$
is in any case positive in one of the two factors. Equation 4.
yields, if one takes on the right the residue with respect to the modulus
145,
\weq{4}{\delta^2 \equiv -1 \pmod{145},}
and this congruence has the four solutions
\[
\delta \equiv \pm 12,\ \pm 17 \pmod{145}.
\]
Since $\delta$ must moreover be even, by 4., $\delta$ could have the following
values:
\[
\delta = \pm 12,\ \pm 128,\ \pm 162.
\]
For the time being one will go no further, since large numerical values for
$\delta$ are improbable from the outset. Corresponding to these values of $\delta$
one obtains from 4. for $\gamma$:
\[
\gamma = -49,\ 7,\ 41.
\]
If we apply in addition the modulus 5, it follows that:
\[
\gamma^2 \equiv 1,\ -1,\ 1 \pmod{5}
\]
and from 3.:
\[
\beta \equiv \pm 1,\ \pm 2,\ \pm 1 \pmod{5}
\]
and from 1.:
\[
\alpha^2 \equiv 1 \pm 2,\ 1 \mp 1,\ 1 \pm 2 \pmod{5}.
\]
From this it follows, since only 0, 1 and 4, not 2 and 3, can be
congruent to a square with respect to the modulus 5, that
the signs in $\delta$ must be taken thus:
\[
\begin{aligned}
\delta &= -12,\ +128,\ -162\\
\alpha &\equiv \pm 2,\ 0,\ \pm 2 \pmod{5}.
\end{aligned}
\]
Equation 2. shows further that the middle case is to be excluded, since the right-hand side is not divisible by 5.
